\chapter{Revisão da Literatura}
\label{cap:02}

% =====================================================================
% ROTEIRO DO CAPÍTULO 2
% Cada seção traz, em comentários, o que escrever:
%   ABERTURA      frase sugerida para começar a seção
%   IDEIA         o que o autor afirma e o que usar dele
%   COMO ABORDAR  o ângulo do parágrafo
%   FONTE         \cite{chave}, onde está no livro/artigo e o link
%   PONTE         ideia ou frase que leva ao passo seguinte
% Os itens marcados CONFERIR precisam ser checados na fonte ao escrever.
% Não repetir o Cap. 5, que aplica requisitos, casos de uso e modelo de
% domínio; aqui ficam o conceito geral e o porquê.
% =====================================================================

% ABERTURA DO CAPÍTULO (1 parágrafo curto)
% COMO ABORDAR: apresentar o caminho como um funil, da prática humana
%   (colecionar) às áreas técnicas usadas para construir a plataforma.
% FRASE SUGERIDA: "Este capítulo reúne os fundamentos que sustentam a
%   plataforma, partindo da prática de colecionar e de sua passagem para
%   o meio digital e chegando às áreas da engenharia de software e da
%   interação humano-computador envolvidas na sua construção."
% Fechar com uma frase que liste as sete seções na ordem.


\section{Coleções e Colecionismo Digital}
\label{sec:colecionismo}

% PAPEL DA SEÇÃO: mostrar que colecionar é um fenômeno estudado e que, no
%   digital, as necessidades de organizar e exibir continuam.
% ABERTURA: "Colecionar vai além de acumular objetos, pois envolve escolher
%   o que entra no conjunto, organizá-lo e atribuir a ele um significado
%   pessoal."
%
% PASSO 1. O que é colecionar
% IDEIA: Belk define colecionar como adquirir e possuir de forma ativa,
%   seletiva e apaixonada coisas retiradas do uso comum e vistas como
%   parte de um conjunto de objetos ou experiências não idênticos.
%   "Seletiva" separa coleção de acúmulo. "Experiências" mostra que não
%   precisa ser objeto físico, então filmes vistos e livros lidos contam.
% COMO ABORDAR: citar a definição (pode ser citação direta curta) e
%   destacar esses dois termos com o exemplo de quem coleciona leituras ou
%   filmes assistidos.
% FONTE: \cite{Belk1995}, cap. 3. CONFERIR a página da definição (por
%   volta da p. 67).
%   https://www.routledge.com/Collecting-in-a-Consumer-Society/Belk/p/book/9780415105347
% PONTE: "Se a coleção pode ser feita de experiências, ela também pode
%   existir em meio digital."
%
% PASSO 2. A coleção no meio digital
% IDEIA: Watkins, Sellen e Lindley identificaram três práticas centrais
%   dos colecionadores, a aquisição, a curadoria e a exibição. O ambiente
%   digital pode favorecer ou atrapalhar cada uma delas.
% COMO ABORDAR: apresentar as três práticas uma a uma, com um exemplo de
%   cada (adicionar um item, organizar e anotar, mostrar a alguém). Elas
%   viram o fio condutor do capítulo, pois a curadoria aparece na
%   personalização e a exibição nas redes sociais.
% FONTE: \cite{Watkins2015}.
%   https://dl.acm.org/doi/abs/10.1145/2702123.2702380
% PONTE: "Manter uma coleção digital, porém, não produz automaticamente o
%   mesmo vínculo que a coleção física."
%
% PASSO 3. Posse de coisas digitais
% IDEIA: em entrevistas, Odom e colegas observaram que as pessoas sentem
%   menos posse sobre o que fica guardado em serviços online, e que dar
%   forma, organizar e apresentar esses itens ajuda a criar vínculo.
% COMO ABORDAR: usar como argumento de que a forma de apresentar importa.
%   Não basta guardar o registro, é preciso que ele pareça da pessoa.
%   CONFERIR no artigo as conclusões exatas antes de afirmar.
% FONTE: \cite{Odom2012}. https://doi.org/10.1145/2207676.2207789
% PONTE: "Além do vínculo, o registro tem outra função, que é permitir
%   revisitar a própria trajetória."
%
% PASSO 4. O registro como memória e reflexão
% IDEIA: Li, Dey e Forlizzi propõem um modelo de cinco estágios para
%   sistemas de registro pessoal (preparação, coleta, integração,
%   reflexão e ação). A reflexão depende de a pessoa conseguir rever o
%   que registrou.
% COMO ABORDAR: analogia. Registrar a data em que viu um filme e a nota
%   dada é a coleta, e voltar a esses registros anos depois é a reflexão.
%   Isso dá base teórica ao que a Justificativa conta.
% FONTE: \cite{Li2010}. https://doi.org/10.1145/1753326.1753409
%   PDF aberto: https://courses.cs.washington.edu/courses/cse440/15au/readings/PersonalInformatics-Li2010.pdf
% PONTE / FECHAMENTO: "Construir um sistema que apoie essas práticas
%   exige, antes de tudo, definir com precisão o que ele deve fazer."


\section{Engenharia de Requisitos}
\label{sec:engenharia-requisitos}

% PAPEL DA SEÇÃO: base conceitual curta; o detalhamento fica no Cap. 5.
% ABERTURA: "Todo sistema de software parte de uma definição do que ele
%   deve fazer e das restrições sob as quais deve operar, e essa definição
%   é objeto da engenharia de requisitos."
%
% PASSO 1. Conceito e atividades
% IDEIA: Sommerville trata requisitos como descrições dos serviços do
%   sistema e de suas restrições, e a engenharia de requisitos como o
%   processo de descobrir, analisar, documentar e validar esses
%   requisitos. Pressman e Maxim dividem o processo em tarefas
%   (concepção, levantamento, elaboração, negociação, especificação,
%   validação e gestão).
% COMO ABORDAR: comparar os dois autores. Descrevem o mesmo processo, um
%   de forma mais enxuta e o outro com mais etapas. Mostrar a
%   convergência sem listar tudo em detalhe.
% FONTE: \cite{Sommerville2011}, cap. 4; \cite{pressman2016}, cap. 8,
%   seção 8.1.
% PONTE: "Entre os requisitos levantados, distinguem-se dois tipos."
%
% PASSO 2. Requisitos funcionais e não funcionais
% IDEIA: os funcionais dizem o que o sistema faz. Os não funcionais dizem
%   como ele se comporta (desempenho, segurança, usabilidade) e costumam
%   afetar o sistema inteiro, e não uma função isolada.
% COMO ABORDAR: definir os dois com um exemplo genérico de cada
%   (cadastrar um item; responder em poucos segundos). Encerrar dizendo
%   que a classificação é aplicada no Cap. 5, sem repetir aquele texto.
% FONTE: \cite{Sommerville2011}, seção 4.1.
% PONTE: "Para descrever como o usuário interage com essas funções,
%   usam-se casos de uso."
%
% PASSO 3. Casos de uso
% IDEIA: segundo Valente, o caso de uso descreve a sequência de
%   interações entre um ator e o sistema para alcançar um objetivo.
% COMO ABORDAR: explicar por que a técnica é útil, já que descreve o
%   sistema do ponto de vista de quem usa.
% FONTE: \cite{Valente2020}, cap. 3.
% PONTE / FECHAMENTO: "Definido o que o sistema faz, é preciso
%   representar os conceitos sobre os quais ele opera."


\section{Modelagem Conceitual com UML}
\label{sec:modelagem-conceitual}

% ABERTURA: "Depois de definidos os requisitos, é preciso identificar os
%   conceitos do domínio e as relações entre eles, e a UML oferece uma
%   notação padronizada para isso."
%
% PASSO 1. O que é a UML
% IDEIA: a OMG mantém a UML como linguagem padrão para especificar,
%   visualizar e documentar sistemas. Fowler distingue a perspectiva
%   conceitual (representar o domínio) da perspectiva de software
%   (representar classes do código).
% COMO ABORDAR: apresentar a UML brevemente e usar a distinção de Fowler
%   como gancho, deixando claro que o trabalho usa a perspectiva
%   conceitual.
% FONTE: \cite{OMG2017uml} https://www.omg.org/spec/UML/2.5.1 ;
%   \cite{Fowler2005uml}, cap. 1.
% PONTE: "Na perspectiva conceitual, o artefato central é o modelo de
%   domínio."
%
% PASSO 2. Modelo de domínio
% IDEIA: para Larman, o modelo de domínio é uma representação visual das
%   classes conceituais do mundo real, com atributos e associações.
%   Funciona como um dicionário visual do domínio e não representa
%   objetos de software.
% COMO ABORDAR: definir e reforçar o "não é software". Isso justifica por
%   que o modelo vem antes das decisões de banco e de código.
% FONTE: \cite{Larman2007}, cap. 9.
% PONTE: "Quando o domínio é extenso, o diagrama precisa ser organizado."
%
% PASSO 3. Pacotes
% IDEIA: agrupar classes em pacotes reduz a complexidade visual e mostra
%   quais conceitos são mais próximos entre si.
% COMO ABORDAR: recurso de organização, em uma ou duas frases.
% FONTE: \cite{Larman2007} e \cite{Fowler2005uml}, capítulo de diagramas
%   de pacotes. CONFERIR o número do capítulo em cada livro.
% PONTE / FECHAMENTO: "O modelo conceitual orienta as decisões de
%   arquitetura e de persistência discutidas a seguir."


\section{Arquitetura de Aplicações \textit{Web}}
\label{sec:arquitetura-web}

% ABERTURA: "Aplicações web costumam ser organizadas em partes com
%   responsabilidades distintas, o que facilita compreender, testar e
%   evoluir cada uma delas."
%
% PASSO 1. Camadas
% IDEIA: Fowler descreve três camadas principais (apresentação, domínio e
%   fonte de dados), em que cada camada usa apenas a de baixo. Isso
%   permite entender uma camada isoladamente e trocar sua implementação,
%   ao custo de mudanças que às vezes atravessam várias camadas.
% COMO ABORDAR: apresentar as vantagens e o custo. Mostrar o custo deixa
%   a revisão mais madura.
% FONTE: \cite{Fowler2002}, cap. 1.
% PONTE: "A comunicação entre a camada de apresentação e as demais segue,
%   em muitas aplicações, o estilo REST."
%
% PASSO 2. REST
% IDEIA: Fielding define REST como um conjunto de restrições
%   (cliente-servidor, ausência de estado, cache, interface uniforme e
%   sistema em camadas). Os recursos são identificados por URIs e
%   manipulados por meio de representações.
% COMO ABORDAR: explicar duas restrições com calma. A ausência de estado
%   significa que cada requisição carrega tudo o que é necessário. A
%   interface uniforme se traduz nos métodos e códigos de status do HTTP.
% FONTE: \cite{Fielding2000}, cap. 5; métodos e códigos de status em
%   \cite{fielding2022http}.
% PONTE: "Como cada requisição precisa identificar quem a faz, é comum
%   usar tokens."
%
% PASSO 3 (OPCIONAL). Token de autenticação
% IDEIA: o JWT é um formato compacto e assinado que leva informações sobre
%   o usuário. O servidor confere a assinatura sem precisar guardar a
%   sessão em memória.
% COMO ABORDAR: apresentar como consequência natural da ausência de
%   estado do REST.
% FONTE: \cite{Jones2015jwt}. https://www.rfc-editor.org/rfc/rfc7519
%
% PASSO 4. Single-Page Application
% IDEIA: na SPA a aplicação é carregada uma vez e o navegador passa a
%   atualizar a tela sem recarregar a página, o que aproxima a
%   experiência da de um programa instalado.
% COMO ABORDAR: contrastar com o site tradicional, em que cada clique
%   pede uma página nova ao servidor. Mostrar que isso importa para telas
%   muito interativas.
% FONTE: \cite{Mikowski2013}, cap. 1.
%   https://www.manning.com/books/single-page-web-applications
% PONTE: "Do outro lado, os dados precisam ser guardados de forma
%   consistente."
%
% PASSO 5. Modelo relacional
% IDEIA: Codd propõe representar os dados como relações (tabelas),
%   independentes da forma física de armazenamento. Elmasri e Navathe
%   explicam chaves e restrições de integridade (de entidade e
%   referencial).
% COMO ABORDAR: partir do artigo original como marco histórico e usar o
%   livro para chave primária, chave estrangeira e integridade
%   referencial.
% FONTE: \cite{Codd1970} https://doi.org/10.1145/362384.362685 ;
%   \cite{Elmasri2011}, caps. 3 e 5.
% PONTE: "Programas orientados a objetos, no entanto, não pensam em
%   tabelas."
%
% PASSO 6. Descompasso objeto-relacional e ORM
% IDEIA: Ireland e colegas classificam as diferenças entre objetos
%   (identidade, herança, referências) e tabelas (linhas, chaves) como
%   descompasso de impedância. Ferramentas de mapeamento
%   objeto-relacional fazem essa ponte. O padrão Data Mapper de Fowler
%   separa os objetos do código que os grava no banco.
% COMO ABORDAR: problema primeiro, solução depois. Justifica a existência
%   de ferramentas de ORM sem citar a ferramenta usada (ela está no
%   Cap. 3).
% FONTE: \cite{Ireland2009} https://ieeexplore.ieee.org/document/5071809 ;
%   \cite{Fowler2002}, cap. 10.
% PONTE / FECHAMENTO: "Sobre essa base técnica, a parte visível do
%   sistema é a interface, tema das próximas seções."


\section{Personalização de Interfaces}
\label{sec:personalizacao}

% PAPEL DA SEÇÃO: coração teórico do trabalho; pode ser a seção mais
%   longa do capítulo.
% ABERTURA: "Personalizar uma interface é adaptá-la a quem a utiliza,
%   seja no conteúdo exibido, seja na aparência, e essa adaptação pode
%   partir do sistema ou do próprio usuário."
%
% PASSO 1. Conceito e quem personaliza
% IDEIA: Fan e Poole definem personalização como o processo de mudar a
%   função, a interface, o conteúdo ou a aparência de um sistema para
%   aumentar sua relevância pessoal, e classificam as formas de fazer
%   isso conforme quem toma a iniciativa. Sundar e Marathe chamam de
%   personalização a mudança feita pelo sistema e de customização a feita
%   pelo usuário. No experimento deles, usuários mais experientes
%   preferiram customizar.
% COMO ABORDAR: definir com Fan e Poole e refinar com Sundar e Marathe.
%   Deixar claro que o trabalho se situa no lado da customização, em que
%   o usuário decide.
% FONTE: \cite{Fan2006} https://doi.org/10.1080/10919392.2006.9681199 ;
%   \cite{Sundar2010} https://doi.org/10.1111/j.1468-2958.2010.01377.x
% PONTE: "Se a escolha é do usuário, cabe perguntar por que ele escolhe
%   personalizar."
%
% PASSO 2. Por que as pessoas personalizam a aparência
% IDEIA: Blom e Monk construíram uma teoria a partir de estudos com
%   usuários de computadores e celulares. Certas disposições levam a
%   pessoa a personalizar, e a personalização produz efeitos cognitivos
%   (facilidade de uso), sociais (expressão de identidade) e emocionais
%   (familiaridade, sentimento de posse, prazer).
% COMO ABORDAR: apresentar os três tipos de efeito e ligar cada um a
%   coleções. O social é a coleção mostrar quem a mantém. O emocional é o
%   sentimento de posse, que retoma Odom na Seção 2.1.
% FONTE: \cite{Blom2003} https://doi.org/10.1207/S15327051HCI1803_1
% PONTE: "Esse vínculo entre personalização e identidade foi testado de
%   forma experimental."
%
% PASSO 3. Controle e identidade
% IDEIA: Marathe e Sundar mostraram que a customização é motivada tanto
%   pelo desejo de controle quanto pelo de identidade, e que customizar a
%   aparência reforça o sentimento de identidade, que por sua vez aumenta
%   a percepção de controle.
% COMO ABORDAR: aprofundar o que a Justificativa citou em uma frase,
%   explicando aqui o experimento e a conclusão.
% FONTE: \cite{Marathe2011}.
% PONTE: "Deixar a decisão com o usuário também se mostra vantajoso em
%   termos de uso."
%
% PASSO 4. Adaptável e adaptativo
% IDEIA: no experimento de Findlater e McGrenere, a maioria dos
%   participantes preferiu o menu que eles mesmos ajustavam ao menu que o
%   sistema mudava sozinho.
% COMO ABORDAR: evidência empírica a favor do controle pelo usuário, em
%   poucas frases.
% FONTE: \cite{Findlater2004} https://doi.org/10.1145/985692.985704
% PONTE: "Além do controle, a aparência em si influencia a forma como o
%   sistema é percebido."
%
% PASSO 5. Estética e emoção
% IDEIA: Tractinsky, Katz e Ikar encontraram forte relação entre achar
%   uma interface bonita e achá-la fácil de usar. Norman descreve três
%   níveis do design (visceral, comportamental e reflexivo); o reflexivo
%   trata do significado e da autoimagem.
% COMO ABORDAR: Tractinsky mostra que a estética não é detalhe. O nível
%   reflexivo de Norman explica por que uma coleção com a cara do dono
%   importa.
% FONTE: \cite{Tractinsky2000} https://doi.org/10.1016/S0953-5438(00)00031-X ;
%   \cite{Norman2008}.
% PONTE: "A liberdade de personalizar, porém, precisa preservar a
%   legibilidade."
%
% PASSO 6. Liberdade com legibilidade
% IDEIA: a heurística de flexibilidade e eficiência de uso de Nielsen
%   recomenda permitir customização. O critério 1.4.3 da WCAG exige
%   contraste mínimo de 4,5:1 para texto.
% COMO ABORDAR: mostrar o equilíbrio entre deixar o usuário escolher
%   cores e garantir que o resultado continue legível.
% FONTE: \cite{Nielsen1994} https://www.nngroup.com/articles/ten-usability-heuristics/ ;
%   \cite{campbell2024wcag}.
% PONTE / FECHAMENTO: retomar o Problema do Cap. 1 (ferramentas de campos
%   fixos não deixam decidir a aparência) e dizer que personalizar a
%   aparência das fichas depende de uma forma de edição acessível ao
%   usuário comum, tema da próxima seção.


\section{Manipulação Direta e Editores Visuais}
\label{sec:manipulacao-direta}

% ABERTURA: "Uma forma de dar ao usuário comum o poder de definir a
%   aparência de algo é permitir que ele aja diretamente sobre a
%   representação na tela, sem precisar de comandos ou código."
%
% PASSO 1. O conceito
% IDEIA: Shneiderman reúne sob o nome de manipulação direta as interfaces
%   com três traços. Os objetos de interesse ficam sempre visíveis. As
%   ações são físicas, como arrastar e clicar, em vez de comandos
%   digitados. As operações são rápidas, reversíveis e têm efeito visível
%   na hora.
% COMO ABORDAR: apresentar os três princípios com um exemplo de cada,
%   arrastando e redimensionando um elemento numa folha.
% FONTE: \cite{Shneiderman1983} https://doi.org/10.1109/MC.1983.1654471
%   PDF aberto: https://www.cl.cam.ac.uk/~pr10/iui/shneiderman83.pdf
% PONTE: "Hutchins, Hollan e Norman procuraram explicar por que essas
%   interfaces parecem tão naturais."
%
% PASSO 2. Por que funciona
% IDEIA: a sensação de interação direta vem de duas coisas. Uma é a
%   distância pequena entre o que a pessoa quer fazer e o que precisa
%   fazer (golfos de execução e de avaliação). A outra é a sensação de
%   engajamento, de estar mexendo no próprio objeto.
% COMO ABORDAR: explicar os dois golfos com um exemplo simples (quero a
%   imagem maior, puxo o canto, vejo ela crescer).
% FONTE: \cite{Hutchins1985} https://doi.org/10.1207/s15327051hci0104_2
% PONTE: "Esse estilo de interação se insere em uma classificação mais
%   ampla."
%
% PASSO 3. Tipos de interação
% IDEIA: Rogers, Sharp e Preece classificam a interação em instruir,
%   conversar, manipular e explorar. Manipular é agir sobre objetos
%   virtuais como se faria com objetos físicos.
% COMO ABORDAR: situar o editor de layouts no tipo "manipular". Para uma
%   fonte brasileira, usar Barbosa et al.
% FONTE: \cite{Rogers2013}; ou \cite{Barbosa2021} https://leanpub.com/ihc-ux
% PONTE: "Como a manipulação direta incentiva a experimentação, ela pede
%   um meio de voltar atrás."
%
% PASSO 4. Poder desfazer
% IDEIA: a heurística de controle e liberdade do usuário pede uma "saída
%   de emergência" para ações feitas por engano.
% COMO ABORDAR: uma ou duas frases ligando reversibilidade (Shneiderman)
%   e controle (Nielsen).
% FONTE: \cite{Nielsen1994}.
% PONTE / FECHAMENTO: editores no estilo do Canva aplicam esses
%   princípios, e montar o layout de uma ficha é uma tarefa de
%   manipulação direta.


\section{Redes Sociais e Compartilhamento de Coleções}
\label{sec:redes-sociais}

% ABERTURA: "Exibir a coleção a outras pessoas faz parte do próprio ato de
%   colecionar, e a internet ampliou essa exibição por meio das redes
%   sociais."
%
% PASSO 1. Definição clássica
% IDEIA: boyd e Ellison definem site de rede social como um serviço que
%   permite criar um perfil, montar uma lista de conexões e percorrer as
%   conexões de outros usuários.
% COMO ABORDAR: apresentar os três elementos e notar que o centro é a
%   pessoa e suas conexões.
% FONTE: \cite{boyd2007} https://doi.org/10.1111/j.1083-6101.2007.00393.x
% PONTE: "Com o tempo, o centro dessas plataformas deixou de ser o perfil
%   e passou a ser o fluxo de atualizações."
%
% PASSO 2. Definição revisada com o fluxo de conteúdo
% IDEIA: em 2013 as mesmas autoras revisaram a definição, acrescentando a
%   possibilidade de consumir e interagir com fluxos de conteúdo
%   produzidos pelas conexões, ou seja, o feed.
% COMO ABORDAR: comparar as duas definições e destacar a mudança; o feed
%   passa a ser elemento central. CONFERIR a página da definição no PDF.
% FONTE: \cite{Ellison2013} https://doi.org/10.1093/oxfordhb/9780199589074.013.0008
%   Pré-print: https://public.websites.umich.edu/~enicole/Ellisonboyd_sociality2013.pdf
% PONTE: "Na literatura brasileira, Recuero analisa essas redes a partir
%   dos atores e dos laços entre eles."
%
% PASSO 3. Atores e laços
% IDEIA: Recuero descreve a rede social como atores (pessoas) ligados por
%   conexões (laços), que podem ser mais fortes (amizade, interação
%   frequente) ou mais fracos (acompanhar alguém).
% COMO ABORDAR: usar os laços de intensidades diferentes para mostrar que
%   há formas distintas de se ligar a alguém, como ser amigo ou só
%   acompanhar uma coleção.
% FONTE: \cite{Recuero2009}. CONFERIR os capítulos sobre laços.
% PONTE: "Plataformas que combinam acervo pessoal e rede social já foram
%   objeto de estudo."
%
% PASSO 4. Exemplo estudado
% IDEIA: Thelwall e Kousha analisaram usuários do Goodreads e concluíram
%   que o site é um híbrido, nem só um catálogo de livros nem só uma rede
%   social.
% COMO ABORDAR: evidência de que juntar registro de coleção e interação
%   social é um modelo reconhecido. Não comparar produtos aqui, porque
%   isso é do Cap. 11.
% FONTE: \cite{Thelwall2017} https://doi.org/10.1002/asi.23733
% FECHAMENTO DO CAPÍTULO: retomar as três práticas de Watkins (aquisição,
%   curadoria e exibição) e mostrar que as seções cobriram a base para
%   cada uma.
